\section{Decision Tree Results}

This section summarizes the performance of the Decision Tree classifier using both a scratch implementation and the scikit-learn library. Three resampling strategies were evaluated: SMOTE, Tomek, and SMOTE-Tomek.  

\subsection{Scratch Implementation}

\begin{table}[H]
\centering
\caption{Scratch Decision Tree Performance Across Resampling Strategies}
\begin{tabular}{|l|c|c|c|c|}
\hline
\textbf{Resampling} & \textbf{Training Acc.} & \textbf{Validation Acc.} & \textbf{Test Acc.} & \textbf{Notes on Class 1} \\ \hline
SMOTE & 0.8473 & 0.7958 & 0.7970 & Low precision (0.20), high recall (0.84) \\ \hline
Tomek & 0.9558 & 0.9408 & 0.9500 & Moderate precision (0.62), low recall (0.36) \\ \hline
SMOTE-Tomek & - & - & - & Balanced classes (15,133 per class), improved minority representation \\ \hline
\end{tabular}
\end{table}

\subsection{Scikit-Learn Implementation}

\begin{table}[H]
\centering
\caption{Scikit-Learn Decision Tree Performance Across Resampling Strategies}
\begin{tabular}{|l|c|c|c|c|c|}
\hline
\textbf{Resampling} & \textbf{Training Time (s)} & \textbf{Training Acc.} & \textbf{Validation Acc.} & \textbf{Class 0 F1} & \textbf{Class 1 F1} \\ \hline
SMOTE & 2.25 & 0.9374 & 0.8749 & 0.93 & 0.39 \\ \hline
Tomek & 0.59 & 0.9681 & 0.9382 & 0.97 & 0.38 \\ \hline
SMOTE-Tomek & 2.20 & 0.9422 & 0.8818 & 0.93 & 0.39 \\ \hline
\end{tabular}
\end{table}

\subsection{Summary}

The tables above show that:  

\begin{itemize}
    \item Both scratch and scikit-learn implementations achieve high accuracy on the majority class (Class 0).
    \item Minority class prediction (Class 1) remains challenging; SMOTE and SMOTE-Tomek generally improve recall for rare events compared to Tomek-only resampling.
    \item Scikit-learn models are significantly faster and slightly more accurate than the scratch implementation.
    \item Resampling strategy has a notable effect on minority class performance, highlighting the importance of handling class imbalance in forest fire datasets.
\end{itemize}



\subsection{Bayesian Optimization for Decision Tree}

To further improve the performance of the Decision Tree on the SMOTE-Tomek dataset, Bayesian Optimization was applied to tune the hyperparameters. The search space for optimization included:

\begin{itemize}
    \item \texttt{max\_depth}: 5 to 25
    \item \texttt{min\_samples\_split}: 5 to 50
    \item \texttt{min\_samples\_leaf}: 1 to 20
\end{itemize}

The optimization process ran for 20 iterations and evaluated combinations of hyperparameters to maximize cross-validation accuracy. Table~\ref{tab:bayes_iterations} shows the best-performing iterations.

\begin{table}[H]
\centering
\caption{Selected Iterations from Bayesian Optimization}
\label{tab:bayes_iterations}
\begin{tabular}{|c|c|c|c|c|}
\hline
\textbf{Iteration} & \textbf{CV Accuracy} & \textbf{max\_depth} & \textbf{min\_samples\_split} & \textbf{min\_samples\_leaf} \\ \hline
2 & 0.9380 & 16.97 & 12.02 & 3.96 \\ 
5 & 0.9364 & 21.65 & 14.56 & 4.45 \\ 
6 & 0.9345 & 21.96 & 6.97 & 6.80 \\ 
10 & 0.9365 & 24.14 & 26.29 & 1.0 \\ 
12 & 0.9279 & 25.0 & 25.42 & 10.77 \\ 
14 & 0.9299 & 25.0 & 37.48 & 1.0 \\ 
\hline
\end{tabular}
\end{table}

The Bayesian Optimization process completed in approximately 108.5 seconds. The optimal hyperparameters identified were:

\begin{itemize}
    \item \texttt{max\_depth}: 17
    \item \texttt{min\_samples\_split}: 12
    \item \texttt{min\_samples\_leaf}: 4
    \item \texttt{criterion}: gini
    \item \texttt{random\_state}: 42
\end{itemize}

Using these optimized parameters, the final Decision Tree trained on the SMOTE-Tomek dataset achieved a validation accuracy of \textbf{0.9137}, representing an improvement over default hyperparameter settings.
